\scanpage{219}
\mspage{112}

comme une structure

\[
\bigl(S,S^\vee,R\subset S\times S^\vee,Z\subset\operatorname{Aut}(\underbrace{S\unsure{,|\times}S^\vee,R}_{\Gamma}),\Pi\subset\operatorname{Bij}(S,S^\vee)\bigr)\tag*{(81)}
\]

où, pour mémoire

\[
\left\{\begin{aligned}
&\text{a) } S,S^\vee\text{ deux ensembles, }R\subset S\times S^\vee\text{ une relation}\\
&\text{b) } Z\text{ un sous-groupe \uline{commutatif} de }\operatorname{Aut}(\Gamma)\\
&\text{c) } \Pi\text{ un ensemble de bijections }S\longrightarrow S^\vee,\text{ satisfaisant}\\
&\text{(1) condition de symétrie \unsure{ci-dessus|ci-dev.|cf.} \unsure{de|du} }Z\ \unsure{\ne|?}\ \unsure{\pi|\Pi|70}\\
&\text{(2) }\lambda_{S^\vee}\varphi=\varphi\lambda_S^{-1},\text{ i.e. }\lambda_{S^\vee}\varphi\lambda_S=\varphi\qquad\forall\lambda\in Z\\
&\deleted{\text{\unsure{i.e.} \unsure{cette condition|cette condit.}}}\\
&[\,\mathrm{NB}\ \text{les }\varphi\text{ satisfaisant à \unsure{ces|aux} conditions (1,2)}\\
&\qquad\text{forment un }\inserted{\text{ens }\Pi_Z\ \text{où }Z\text{ opère \uline{librement}}}\ \deleted{\text{ensemble}}\ \deleted{\unsure{G|g}\subset\operatorname{Aut}(\Gamma)}\\
&\qquad\text{par }T_g\varphi=g_{S^\vee}\varphi=\varphi\circ g_S^{-1}\,]\\
&\text{(3) }\Pi\text{ est un sous-}Z\text{-torseur de }\Pi_Z\\
&\deleted{\text{\unsure{cad} }\forall\varphi\in\Pi\ \text{(\unsure{$Z\ni$} pour un }\unsure{\varphi_0}\in\Pi,\text{ c'est kif-kif)}}
\end{aligned}\right.\tag*{(82)}
\]

% slanted
\margin{La condition (1), \unsure{assure|sous} que les données a) b) c), définissant un \inserted{ens. de} \deleted{parties} structures de \unsure{parties|graphes} \deleted{\unsure{strict}} strict séparés (moins \unsure{exiger} que $R$ soit séparante), ayant \deleted{un} \unsure{pr conséquence} \unsure{sous-entendu}, la condition (2) que $Z$ est un groupe d'homothéties pour cette structure de graphe.}

% slanted
\margin{\uline{NB} Il suffit de poser (2) pour \uline{un} $\varphi_0\in\Pi$, pour}

Si on considère l'objet

\[
\bigl(S'=S^\vee,S'^\vee=S,R'={}^sR,Z'={}^sZ,\Pi'=\Pi^{-1}=\{\varphi^{-1}\mid\varphi\in\Pi\}\bigr),\tag*{(83)}
\]

on voit qu'elle satisfait aux \deleted{\unsure{m}} conditions (82) que le système de départ (80). Notons que la bijection canonique

\[
\varphi\longmapsto\varphi^{-1}\qquad\Pi\xrightarrow{\sim}\Pi'\tag*{(84)}
\]

satisfait

\[
\bigl(\underbrace{T_\lambda\varphi}_{\lambda_{S^\vee}\varphi}\bigr)^{-1}
=\underbrace{T'_{\lambda'}(\varphi^{-1})}_{\deleted{\unsure{\lambda}}\ \varphi^{-1}\lambda_{S'}^{-1}}
\qquad\text{grand succès \unsure{!|?}},\tag*{(85)}
\]

% slanted
\margin{dans cette formule, j'identifie, \unsure{pr|je} \unsure{pos.|pose} $\lambda'=(\lambda_{S^\vee},\lambda_S)$ à $\lambda=(\lambda_S,\lambda_{S^\vee})$}

\margin{pour l'iso $\lambda\longmapsto\lambda'$, $Z\longrightarrow Z'$.}

i.e. (84) est \deleted{\unsure{équiv.}} un iso de $Z$-ensembles ! Si on considère d'autre part l'iso déduit de (76) en \deleted{\unsure{ten}} faisant le twist par \deleted{$\Pi^{\unsure{(-1)|\mathrm{inv}}}$} le torseur inverse de \deleted{$\Pi$}

\[
S!\simeq\Pi^{\mathrm{inv}}\wedge_Z S^\vee
\]

on en déduit, en passant à l'action opposée de $Z$ sur les deux mb

\[
S\simeq\Pi\wedge_Z(S^\vee!),\tag*{(86)}
\]

% archive sheet number: 218
